\originalchapter{搜索与最短路}\label{ch:shortest}
\sourceof{6.042J 第9章有向图与调度部分; 本章算法推导为编者补充}
\originalsection{表示与搜索}
图算法的运行时间依赖输入如何存储. 邻接矩阵是 $n\times n$ 矩阵, 第 $uv$ 项记录是否有边或其权. 它可在常数时间查询一条边, 但占 $O(n^2)$ 空间. 邻接表为每点列出邻边, 总大小 $O(n+m)$; 搜索稀疏图通常采用邻接表. 多重边需分配独立边编号, 不能只用端点识别.

广度优先搜索 BFS 从起点 $s$ 出发, 使用先进先出队列. 初始化 $d(s)=0$, 其他距离为未知. 取出队首 $u$, 扫描其每个邻点 $v$; 若 $v$ 未访问, 设置 $d(v)=d(u)+1$, 父顶点为 $u$, 并入队. 每点最多入队一次, 每条无向边扫描两次, 有向弧扫描一次, 总时间 $O(n+m)$.
\begin{theorem}\label{thm:bfs}
BFS 首次访问顶点时记录的距离等于无权图最短路长度. 父边在可达部分构成最短路树.
\end{theorem}
\begin{proof}
队列按记录距离非递减取出: 距离 $k$ 点新加入的邻点全为 $k+1$, 它们排在尚存距离 $k$ 点之后. 父边给出一条长度 $d(v)$ 的实际路, 所以真实距离不大于记录值. 反向按真实距离归纳, 若最短路中 $v$ 的前点 $u$ 距离为 $k-1$, 它先被取出时将 $v$ 至迟以距离 $k$ 访问. 因而记录值不大于真实距离. 父顶点距离严格减1, 所以无圈, 沿父边都回到根.
\end{proof}
深度优先搜索 DFS 则沿尚未访问的邻点递归, 无新邻点时回退. 顶点分为白色未访问、灰色在递归栈中、黑色已完成. 每点进入和退出各一次, 也需 $O(n+m)$ 时间. 在有向图中, 指向灰色顶点的弧与递归栈上的路组成有向圈; 有向圈在 DFS 中必产生这样的弧, 因为从圈上最先访问点继续追踪时, 在它退出之前圈上可达点都必须被处理.
\originalsection{有向无圈图与调度}
\begin{definition}
有向无圈图简称 DAG. 拓扑排序是顶点的线性顺序, 使每条弧都由较前顶点指向较后顶点.
\end{definition}
\begin{theorem}\label{thm:toposort}
有限有向图有拓扑排序当且仅当它是 DAG. 反复删除入度0顶点即可构造排序.
\end{theorem}
\begin{proof}
拓扑顺序沿弧严格前进, 不可能形成有向圈. 若 DAG 没有入度0顶点, 从任一点不断沿入弧倒走, 有限顶点保证重复, 产生有向圈, 矛盾. 删除一个入度0点后仍无圈, 归纳排序余图, 把删点放在最前即可. 用队列维护入度, 每删除点时减少其后继的入度, 每弧处理一次, 时间 $O(n+m)$.
\end{proof}
若每项工作 $v$ 耗时 $p(v)\ge0$, 弧 $u\to v$ 表示 $v$ 必须等 $u$ 完成, 无限并行资源下最早完工时间满足 $F(v)=p(v)+\max_{u\to v}F(u)$, 没有前驱时最大值记0. 按拓扑顺序计算, 每个前驱值都已确定. 任意调度的完成时间至少达到这个值, 因最长先后依赖链必须串行; 按这些最早开始时间安排又满足全部约束, 所以这是可实现的最优工期. 资源人数限制会改变问题, 不能直接沿用此结论.
\originalsection{松弛与 Dijkstra}
对弧 $uv$ 赋实权 $w(uv)$, 路权为边权和. 若起点可达负权圈, 经过它重复行走可使某些终点路权趋于 $-\infty$, 此时不存在有限最短行走. Dijkstra 假设全部边权非负.

距离松弛是把估计值 $D(v)$ 改为 $\min\{D(v),D(u)+w(uv)\}$. 从 $D(s)=0$, 其余为 $+\infty$ 出发, 每个有限估计都对应实际行走, 因而它始终是最短距离的上界. Dijkstra 反复从未确定点中选最小估计点 $u$, 将它确定, 再松弛出弧.
\begin{theorem}\label{thm:dijkstra}
非负权图中, Dijkstra 确定 $u$ 时有 $D(u)=\dist_w(s,u)$.
\end{theorem}
\begin{proof}
假设已有确定点的值均正确. 对新选 $u$, 取一条最短路; 非负权保证可删圈而得到简单最短路. 若 $u$ 不可达, 所有未确定点估计均无限, 算法可以结束. 否则令 $y$ 为该路上第一个未确定点, 前点 $x$ 已确定, 松弛后 $D(y)\le\dist_w(s,y)$. 结合上界不变量, 等号成立. 非负性保证路径后段权非负, 所以 $\dist_w(s,y)\le\dist_w(s,u)$. 选择最小估计给 $D(u)\le D(y)\le\dist_w(s,u)$, 与始终为上界合起来得等号. 起点为归纳基点.
\end{proof}
用矩阵或顺序扫描选最小值, 时间 $O(n^2+m)$. 用邻接表与带顶点位置索引的二叉堆, 每点在堆中至多一个项, 成功松弛做减小键值, 时间 $O((n+m)\log n)$. 若改为加入新项并跳过过期项, 堆最多有 $O(n+m)$ 项, 对任意多重图应写 $O((n+m)\log(n+m))$; 在简单图 $m\le n(n-1)$ 时仍可写前一界. 算法需要的非负条件是正确性条件, 不只是程序优化建议.
\begin{example}
有向弧为 $sa:4$, $sb:1$, $ba:2$, $ac:1$, $bc:5$, $ct:3$, $at:7$. 从 $s$ 求最短距离.
\end{example}
\begin{solution}
先确定 $s$, 得 $D(a)=4,D(b)=1$. 确定 $b$ 后, $a$ 改为3, $c$ 为6. 确定 $a$ 后, $c$ 改为4, $t$ 为10. 确定 $c$ 后, $t$ 改为7. 最终顺序 $s,b,a,c,t$, 距离 $0,1,3,4,7$. 最短路 $s,b,a,c,t$ 的权为 $1+2+1+3=7$, 父边同时提供可验证路径.
\end{solution}
\originalsection{负权与 Bellman--Ford}
\begin{theorem}\label{thm:bellmanford}
定义 $D_k(v)$ 为从 $s$ 到 $v$ 使用至多 $k$ 条弧的最小行走权, 不存在时为无限. 则
\[
D_{k+1}(v)=\min\bigl\{D_k(v),\min_{uv\in E}(D_k(u)+w(uv))\bigr\}.
\]
若没有从 $s$ 可达的负圈, $D_{n-1}$ 为最短距离. 若 $D_n(v)<D_{n-1}(v)$ 对某点成立, 则有可达负圈.
\end{theorem}
\begin{proof}
至多 $k+1$ 弧的行走或者原已至多 $k$ 弧, 或最后一步为 $uv$, 之前至多 $k$ 弧, 得递推. 无可达负圈时可删除行走中的非负圈, 最短值由至多 $n-1$ 弧的简单路实现. 若第 $n$ 轮仍严格改进, 一条改进的行走有 $n$ 条弧而至少重复一点. 若所含所有圈非负, 删圈得到至多 $n-1$ 弧且权不增加的行走, 与严格改进矛盾; 因而含可达负圈.
\end{proof}
上式用两份数组保存相邻轮, 是最容易核查的不变量版本, 时间 $O(nm+n^2)$, 空间 $O(n)$ 可再配父指针. 通常把 $n-1$ 轮所有边松弛称为 Bellman--Ford; 原地更新可以在一轮传播多条边, 但仍保证最终结果, 因为它不会增加估计且至少包含上述一轮的改进. 第 $n$ 轮能够改进的点及其后继, 最短行走值为 $-\infty$; 与负圈无关的点仍可能有有限最短距离.
\originalsection{全源最短路}
Floyd--Warshall 令 $D^{(k)}_{ij}$ 为内部顶点只允许来自 $\{1,\ldots,k\}$ 的最短路权. 非对角初值为直接弧最小权, 无弧时为无限; 对角初值为0与该点所有自环权的最小值, 没有自环时为0. 这样不会漏掉负权自环. 在无负圈条件下,
\[
D^{(k)}_{ij}=\min\{D^{(k-1)}_{ij},D^{(k-1)}_{ik}+D^{(k-1)}_{kj}\}.
\]
一条简单最短路或不经过 $k$, 或在 $k$ 分成两段, 所以递推正确. 三重循环时间 $O(n^3)$, 矩阵空间 $O(n^2)$. 若最终某个对角值为负, 图有负圈; 必须报告这个条件, 不能把负对角仍称为正常有限最短路矩阵.
\originalsection{习题与解答}
\begin{exercise}\label{ex:short1}
构造一张有负权弧而无负圈的图, 使 Dijkstra 的确定步骤出错.
\end{exercise}
\begin{exercise}\label{ex:short2}
无向图每边权为1. 解释为何 BFS 比二叉堆 Dijkstra 更直接, 并说明权只取0、1时可怎样用双端队列.
\end{exercise}
习题\ref{ex:short1}的解答.
\begin{solution}
取弧 $sa:2,sb:5,ba:-4$, 图无有向圈. Dijkstra 先确定 $a$ 为2, 但真正距离沿 $s,b,a$ 为1. 负权使最短路上较后点距离反而更小, 破坏证明中的非递减关系.
\end{solution}
习题\ref{ex:short2}的解答.
\begin{solution}
单位权下新点距离恰比出队点多1, 普通队列自动保持层序, 时间 $O(n+m)$. 0、1权下每次成功松弛权0边将邻点放队首, 权1边放队尾, 保存过期距离用于跳过旧项, 就维持队列按距离层次排列. 从最小层向外处理的正确性与非负权证明一致, 这是0--1 BFS; 用标准的确定标志实现时每点仅在首次以最小距离出队时扫描邻边, 时间仍为 $O(n+m)$.
\end{solution}
